\documentclass[10pt,a4paper]{amsart}
\usepackage[margin=19mm]{geometry}
\usepackage{amsmath,amssymb,mathtools}
\usepackage[scr=rsfs,cal=euler]{mathalfa}
\usepackage{xfrac}
\usepackage[unicode,hidelinks,bookmarks=false]{hyperref}
\newcommand{\ie}{i.e.\ }
\newcommand{\cf}{cf.\ }
\newcommand{\resp}{respectively\ }
\newcommand{\Subsection}{\subsection}
\providecommand{\nobreakdash}{\nobreak\hbox{-}}
\setlength{\emergencystretch}{2em}
\multlinegap0pt
\usepackage{amssymb}
\usepackage[normalem]{ulem}
\usepackage[shortlabels]{enumitem}
\newtheorem{lemma}{Lemma}
\newcommand{\OO}{{\mathcal O}}
\newcommand{\Ome}{\Omega}
\renewcommand{\d}{\delta}
\newcommand{\ddt}{\, \frac{d}{dt}}
\newcommand{\upref}[2]{\hspace{-0.8ex}\stackrel{\eqref{#1}}{#2}} % standard
\title{Solutions of the thin film equation obtained in the limit of vanishing slip}
\author{Hans Knüpfer, Juan J.~L. Velázquez}
\date{Source: arXiv:2512.17463v1 (CC BY 4.0)}
\begin{document}
\maketitle

\noindent\emph{Source and licence.} Solutions of the thin film equation obtained in the limit of vanishing slip. Version arXiv:2512.17463v1 (CC BY 4.0), distributed by arXiv under the Creative Commons Attribution 4.0 International licence, http://creativecommons.org/licenses/by/4.0/, as recorded in the arXiv licence metadata for this exact version and shown on the abstract page https://arxiv.org/abs/2512.17463v1. Full licence notice: \texttt{licenses/arxiv-2512.17463-CC-BY-4.0.txt}.\par\smallskip

\noindent\textit{Editorial scope.} Counted unit: the article's lemma lemma (source0.tex lines 1372--1390). The article's complete original proof of it is delivered with it (source0.tex lines 1391--1426, one \texttt{proof} environment of 36 lines). No mathematics is written, repaired or re-derived: the statement and the proof are the article's own lines, in the article's own notation, and no retained line is substituted or re-flowed.  Retained uncounted as the source-local context the proof needs: lines 762--765.  Nothing was omitted from any retained span, so the package carries no omission record.  No retained line contains a citation, so the wrapper carries no bibliography and no bibliography entry.  No retained line contains a figure environment, an image-inclusion command or drawing code, so no image is delivered and the package declares no asset dependency.  Editorial formatting only, in addition: an AMS \texttt{amsart} preamble that restates the article's own declarations and macros used by the retained lines, namely the environments \texttt{lemma} on separate counters, together with the standard packages those lines need.  The retained environments print in this document as Lemma 1 in order of appearance, whereas the article prints the same items with the numbering of its own full document.  The complete native source of the article as distributed in the arXiv e-print for this exact version is bundled unchanged as \texttt{sources/191315/source0.tex}.
\begin{align} \label{tfe-noslip} 
    h_t +(h^3 h_{xxx})_x \ = \ 0 \qquad\qquad %
  \text{in $\Ome \ := \ \{h>0\}$.}
\end{align}

  \begin{lemma}[Energy dissipation for type (b) solutions]
    Let $h$ be a solution of \eqref{tfe-noslip} with support $(s,\infty)$ for
    some function $s = s(t)$. We assume that $h$ is a type (b) in the sense that
    for $\xi :=x- s$ we have
    \begin{align}
      h(\xi,t) \ %
      \approx \ (3 \dot s)^{\frac 13} \xi |\ln \xi|^{\frac 13} \qquad
      \text{for $\xi \ll 1$.}
    \end{align}
    We also assume that the time derivative $h_t$ and up to three spatial
    derivatives of $h$ have the asymptotic behaviour, obtained by applying these
    derivatives on the right--hand side above.  Furthermore, the solution and its
    derivatives decay sufficiently fast for $\xi\gg 1$. Then the limit below
    exists and is finite and
    \begin{align}
      \frac d{dt} \frac 12   \int_{s}^\infty h_x^2\ dx    \ %
      = \ \lim_{\d \to 0} \Big( \frac{\dot s}2h_x^2(\d )-\int_{s + \d}^{\infty}h^3h_{xxx}^2 \ dx \Big).
    \end{align}
\end{lemma}

\begin{proof}
  By assumption, for $\xi \ll 1$ we have
  $h_{\xi\xi\xi} \approx \frac 13 (3 \dot s)^{\frac 13} \xi^{-2} |\ln \xi|^{-\frac 23}$ and
  \begin{align}  %
    &h_\xi h_t  \ %
    = \ \OO(\xi |\ln \xi |^{\frac23}) \ = \ o(1), \label{asy1} \\ \qquad %
    &h^3h_{\xi\xi}h_{\xi\xi\xi} \ %
    = \ \OO(|\ln \xi |^{-\frac 13}) \ = \ o(1). \label{asy3} %
  \end{align}
  where $o(1) \to 0$ as $\xi \to 0$.  For $\xi > 0$ we have
  $h_t = h_\xi\dot s - (h^3 h_{\xi\xi\xi})_\xi$. Using this identity, for
  $\d > 0$ we hence obtain
  \begin{align}
    \ddt \frac 12 \int_{\d}^{\infty } h_\xi ^2\ dx  \ %
    &\stackrel{p.i.}= \ -\int_{\d}^{\infty }h_{\xi\xi}h_t + \left[ h_{\xi
      }h_t \right]_{\d}^{\infty } \notag \\
    &\upref{asy1}= \ -\dot s\int_{\d}^{\infty } h_{\xi\xi} h_\xi  \ d\xi %
      + \int_{\d}^{\infty }h_{\xi\xi}(h^3h_{\xi\xi \xi})_\xi +  o(1) \notag \\
    &= \ \frac{\dot s}2h_\xi ^2(\d )-\int_{\d}^{\infty} h^3h_{\xi\xi \xi}^2 \ d\xi %
      -\left[ h^3h_{\xi\xi}h_{\xi\xi \xi}  \right]_{\d}^{\infty }+ o(1) \notag \\
    &\upref{asy3}= \ \frac{\dot s}2h_\xi ^2(\d) - \int_{\d}^{\infty}h^3h_{\xi\xi \xi}^2 \ d\xi  + o(1),  \notag
  \end{align}%
  where $o(1) \to 0$ as $\d \to 0$. For $\d \to 0$ we obtain the asserted
  identity. Furthermore, we note that we have  the integral identity
  \begin{align}
    \int_\d^{\frac 12} \frac 1{\xi |\ln \xi|^{\frac 13}} \ %
    \approx \ \frac 32 |\ln \d|^{\frac 23}  \qquad\qquad \text{for $\d \ll 1$}
  \end{align}
  In particular, for  $\xi \to 0$ a straightforward calculation yields
  \begin{align}
    \frac{\dot s}2h_\xi ^2(\d) - \int_{\d}^{\infty}h^3h_{\xi\xi \xi}^2 \ d\xi \ %
    \approx \ \frac 16  (3\dot s)^{\frac 53} |\ln \d|^{\frac 23}  -  \frac 16 (3 \dot s)^{\frac 53} |\ln \d|^{\frac 23} + \OO(1).
  \end{align}
  i.e. the cancellation of the leading order terms shows that the limit in
  leading order exists and is finite.
\end{proof}

\end{document}
